% paper76-copilot-channel.tex — The Copilot Channel:
%   Architecture, Jurisdiction, and Infrastructure
%   Paper 76 of the Judgment Geometry series.
% Compile: pdflatex paper76-copilot-channel
\documentclass[11pt]{article}
\usepackage{lmodern}
% jugeo-common.tex — shared preamble for all Judgment Geometry papers
% Include via: % jugeo-common.tex — shared preamble for all Judgment Geometry papers
% Include via: % jugeo-common.tex — shared preamble for all Judgment Geometry papers
% Include via: \input{jugeo-common}

% ── Packages ────────────────────────────────────────────────────────────
\usepackage{amsmath,amssymb,amsthm,mathtools}
\usepackage{tikz-cd}
\usepackage{listings}
\usepackage{xcolor}
\usepackage{xspace}
\usepackage{booktabs}
\usepackage{multirow}
\usepackage{enumitem}
\usepackage[expansion=false]{microtype}
\usepackage{hyperref}
\usepackage{cleveref}
% \usepackage{thmtools}  % disabled: not available in this TeX installation
\usepackage{float}
\usepackage{caption}
\usepackage{subcaption}
\usepackage{tabularx}
\usepackage{stmaryrd}       % \llbracket, \rrbracket
\usepackage{bbm}            % \mathbbm{1}

% ── Page geometry ───────────────────────────────────────────────────────
\usepackage[margin=1in]{geometry}

% ── Theorem environments ────────────────────────────────────────────────
\theoremstyle{plain}
\newtheorem{theorem}{Theorem}[section]
\newtheorem{lemma}[theorem]{Lemma}
\newtheorem{proposition}[theorem]{Proposition}
\newtheorem{corollary}[theorem]{Corollary}

\theoremstyle{definition}
\newtheorem{definition}[theorem]{Definition}
\newtheorem{example}[theorem]{Example}
\newtheorem{construction}[theorem]{Construction}

\theoremstyle{remark}
\newtheorem{remark}[theorem]{Remark}
\newtheorem{notation}[theorem]{Notation}

% ── Python listings style ──────────────────────────────────────────────
\definecolor{jugeoBlue}{HTML}{2563EB}
\definecolor{jugeoGreen}{HTML}{059669}
\definecolor{jugeoRed}{HTML}{DC2626}
\definecolor{jugeoPurple}{HTML}{7C3AED}
\definecolor{jugeoGray}{HTML}{6B7280}
\definecolor{jugeoBg}{HTML}{F8FAFC}

\lstdefinestyle{jugeo-python}{
  language=Python,
  basicstyle=\ttfamily\small,
  keywordstyle=\color{jugeoBlue}\bfseries,
  stringstyle=\color{jugeoGreen},
  commentstyle=\color{jugeoGray}\itshape,
  numberstyle=\tiny\color{jugeoGray},
  backgroundcolor=\color{jugeoBg},
  numbers=left,
  numbersep=6pt,
  frame=single,
  framerule=0.4pt,
  rulecolor=\color{jugeoGray!40},
  breaklines=true,
  breakatwhitespace=true,
  showstringspaces=false,
  tabsize=4,
  xleftmargin=1.5em,
  framexleftmargin=1.5em,
  morekeywords={dataclass,frozen,field,Enum,IntEnum,auto,Any,
                Mapping,Sequence,Optional,match,case},
  literate={→}{{$\to$}}1 {←}{{$\leftarrow$}}1
           {≤}{{$\leq$}}1 {≥}{{$\geq$}}1 {≠}{{$\neq$}}1
           {∈}{{$\in$}}1 {∀}{{$\forall$}}1 {∃}{{$\exists$}}1
           {⊕}{{$\oplus$}}1 {⊖}{{$\ominus$}}1,
  escapeinside={(*@}{@*)},
}
\lstset{style=jugeo-python}

\lstdefinestyle{jugeo-output}{
  basicstyle=\ttfamily\small,
  backgroundcolor=\color{jugeoBg},
  frame=single,
  framerule=0.4pt,
  rulecolor=\color{jugeoGray!40},
  breaklines=true,
  showstringspaces=false,
  xleftmargin=1.5em,
  framexleftmargin=0.5em,
  numbers=none,
}

% ── JuGeo-specific macros ──────────────────────────────────────────────

% Judgment 8-tuple
\newcommand{\J}{\mathcal{J}}
\newcommand{\Jud}[8]{(#1,\,#2,\,#3,\,#4,\,#5,\,#6,\,#7,\,#8)}
\newcommand{\JudShort}[1]{\J_{#1}}

% Components
\newcommand{\coord}{\mathsf{c}}            % coordinate
\newcommand{\prop}{\varphi}                 % proposition
\newcommand{\carrier}{\mathcal{A}}          % carrier type
\newcommand{\evid}{\mathcal{E}}             % evidence bundle
\newcommand{\oblig}{\mathcal{O}}            % obligations
\newcommand{\obstr}{\mathcal{B}}            % obstructions
\newcommand{\trust}{\mathcal{T}}            % trust annotation
\newcommand{\prov}{\Pi}                     % provenance

% Trust algebra
\newcommand{\Talg}{\mathfrak{T}}            % trust ordered algebra
\newcommand{\Eadm}{\mathcal{E}_{\mathrm{adm}}}
\newcommand{\tleq}{\preceq}                 % trust ordering
\newcommand{\tjoin}{\oplus}                 % conservative join
\newcommand{\tatten}{\ominus}               % attenuation
\newcommand{\tpromote}{\uparrow_\pi}        % promotion
\newcommand{\tdemote}{\downarrow_\chi}       % demotion

% Trust tiers
\newcommand{\Tcontra}{\ensuremath{\mathsf{CONTRADICTED}}}
\newcommand{\Tunver}{\ensuremath{\mathsf{UNVERIFIED}}}
\newcommand{\Tcopilot}{\ensuremath{\mathsf{COPILOT}}}
\newcommand{\Toracle}{\ensuremath{\mathsf{ORACLE}}}
\newcommand{\Truntime}{\ensuremath{\mathsf{RUNTIME}}}
\newcommand{\Tsolver}{\ensuremath{\mathsf{SOLVER}}}
\newcommand{\Tproof}{\ensuremath{\mathsf{PROOF}}}

% Site and geometry
\newcommand{\Site}{\mathcal{S}}             % semantic site
\newcommand{\Cov}{\mathrm{Cov}}            % covering family
\newcommand{\Ob}{\mathrm{Ob}}              % objects
\newcommand{\Mor}{\mathrm{Mor}}            % morphisms
\newcommand{\res}{\mathrm{res}}            % restriction
\newcommand{\Cech}{\check{C}}              % Čech complex
\newcommand{\Hcoh}[1]{H^{#1}}             % cohomology group

% Descent
\newcommand{\descent}{\mathrm{desc}}
\newcommand{\glue}{\mathrm{glue}}
\newcommand{\overlap}[2]{#1 \cap #2}

% Semantic moves
\newcommand{\VERIFY}{\textsc{Verify}}
\newcommand{\CONSTRUCT}{\textsc{Construct}}
\newcommand{\NORMALIZE}{\textsc{Normalize}}
\newcommand{\GLUE}{\textsc{Glue}}
\newcommand{\RESTRICT}{\textsc{Restrict}}
\newcommand{\TRANSPORT}{\textsc{Transport}}
\newcommand{\REPAIR}{\textsc{Repair}}
\newcommand{\PROMOTE}{\textsc{Promote}}
\newcommand{\CHALLENGE}{\textsc{Challenge}}
\newcommand{\ESCALATE}{\textsc{Escalate}}

% Fragments
\newcommand{\QFLIA}{\texttt{QF\_LIA}}
\newcommand{\QFLRA}{\texttt{QF\_LRA}}
\newcommand{\QFBV}{\texttt{QF\_BV}}
\newcommand{\QFUF}{\texttt{QF\_UF}}

% Misc
\newcommand{\Python}{\textsc{Python}}
\newcommand{\JuGeo}{\textsc{JuGeo}}
\newcommand{\LEAN}{\textsc{Lean}\,4}
\newcommand{\Fstar}{F$^{\star}$}
\newcommand{\Coq}{\textsc{Coq}}
\newcommand{\Dafny}{\textsc{Dafny}}

% Common shortcuts
\newcommand{\ie}{i.e.\@\xspace}
\newcommand{\eg}{e.g.\@\xspace}
\newcommand{\cf}{cf.\@\xspace}
\newcommand{\wrt}{w.r.t.\@\xspace}

% Evaluation shortcuts
\newcommand{\numcases}{300}
\newcommand{\numspec}{100}
\newcommand{\numeq}{100}
\newcommand{\numbug}{100}
\newcommand{\medtime}{6.5\,ms}
\newcommand{\totaltime}{1.949\,s}


% ── Packages ────────────────────────────────────────────────────────────
\usepackage{amsmath,amssymb,amsthm,mathtools}
\usepackage{tikz-cd}
\usepackage{listings}
\usepackage{xcolor}
\usepackage{xspace}
\usepackage{booktabs}
\usepackage{multirow}
\usepackage{enumitem}
\usepackage[expansion=false]{microtype}
\usepackage{hyperref}
\usepackage{cleveref}
% \usepackage{thmtools}  % disabled: not available in this TeX installation
\usepackage{float}
\usepackage{caption}
\usepackage{subcaption}
\usepackage{tabularx}
\usepackage{stmaryrd}       % \llbracket, \rrbracket
\usepackage{bbm}            % \mathbbm{1}

% ── Page geometry ───────────────────────────────────────────────────────
\usepackage[margin=1in]{geometry}

% ── Theorem environments ────────────────────────────────────────────────
\theoremstyle{plain}
\newtheorem{theorem}{Theorem}[section]
\newtheorem{lemma}[theorem]{Lemma}
\newtheorem{proposition}[theorem]{Proposition}
\newtheorem{corollary}[theorem]{Corollary}

\theoremstyle{definition}
\newtheorem{definition}[theorem]{Definition}
\newtheorem{example}[theorem]{Example}
\newtheorem{construction}[theorem]{Construction}

\theoremstyle{remark}
\newtheorem{remark}[theorem]{Remark}
\newtheorem{notation}[theorem]{Notation}

% ── Python listings style ──────────────────────────────────────────────
\definecolor{jugeoBlue}{HTML}{2563EB}
\definecolor{jugeoGreen}{HTML}{059669}
\definecolor{jugeoRed}{HTML}{DC2626}
\definecolor{jugeoPurple}{HTML}{7C3AED}
\definecolor{jugeoGray}{HTML}{6B7280}
\definecolor{jugeoBg}{HTML}{F8FAFC}

\lstdefinestyle{jugeo-python}{
  language=Python,
  basicstyle=\ttfamily\small,
  keywordstyle=\color{jugeoBlue}\bfseries,
  stringstyle=\color{jugeoGreen},
  commentstyle=\color{jugeoGray}\itshape,
  numberstyle=\tiny\color{jugeoGray},
  backgroundcolor=\color{jugeoBg},
  numbers=left,
  numbersep=6pt,
  frame=single,
  framerule=0.4pt,
  rulecolor=\color{jugeoGray!40},
  breaklines=true,
  breakatwhitespace=true,
  showstringspaces=false,
  tabsize=4,
  xleftmargin=1.5em,
  framexleftmargin=1.5em,
  morekeywords={dataclass,frozen,field,Enum,IntEnum,auto,Any,
                Mapping,Sequence,Optional,match,case},
  literate={→}{{$\to$}}1 {←}{{$\leftarrow$}}1
           {≤}{{$\leq$}}1 {≥}{{$\geq$}}1 {≠}{{$\neq$}}1
           {∈}{{$\in$}}1 {∀}{{$\forall$}}1 {∃}{{$\exists$}}1
           {⊕}{{$\oplus$}}1 {⊖}{{$\ominus$}}1,
  escapeinside={(*@}{@*)},
}
\lstset{style=jugeo-python}

\lstdefinestyle{jugeo-output}{
  basicstyle=\ttfamily\small,
  backgroundcolor=\color{jugeoBg},
  frame=single,
  framerule=0.4pt,
  rulecolor=\color{jugeoGray!40},
  breaklines=true,
  showstringspaces=false,
  xleftmargin=1.5em,
  framexleftmargin=0.5em,
  numbers=none,
}

% ── JuGeo-specific macros ──────────────────────────────────────────────

% Judgment 8-tuple
\newcommand{\J}{\mathcal{J}}
\newcommand{\Jud}[8]{(#1,\,#2,\,#3,\,#4,\,#5,\,#6,\,#7,\,#8)}
\newcommand{\JudShort}[1]{\J_{#1}}

% Components
\newcommand{\coord}{\mathsf{c}}            % coordinate
\newcommand{\prop}{\varphi}                 % proposition
\newcommand{\carrier}{\mathcal{A}}          % carrier type
\newcommand{\evid}{\mathcal{E}}             % evidence bundle
\newcommand{\oblig}{\mathcal{O}}            % obligations
\newcommand{\obstr}{\mathcal{B}}            % obstructions
\newcommand{\trust}{\mathcal{T}}            % trust annotation
\newcommand{\prov}{\Pi}                     % provenance

% Trust algebra
\newcommand{\Talg}{\mathfrak{T}}            % trust ordered algebra
\newcommand{\Eadm}{\mathcal{E}_{\mathrm{adm}}}
\newcommand{\tleq}{\preceq}                 % trust ordering
\newcommand{\tjoin}{\oplus}                 % conservative join
\newcommand{\tatten}{\ominus}               % attenuation
\newcommand{\tpromote}{\uparrow_\pi}        % promotion
\newcommand{\tdemote}{\downarrow_\chi}       % demotion

% Trust tiers
\newcommand{\Tcontra}{\ensuremath{\mathsf{CONTRADICTED}}}
\newcommand{\Tunver}{\ensuremath{\mathsf{UNVERIFIED}}}
\newcommand{\Tcopilot}{\ensuremath{\mathsf{COPILOT}}}
\newcommand{\Toracle}{\ensuremath{\mathsf{ORACLE}}}
\newcommand{\Truntime}{\ensuremath{\mathsf{RUNTIME}}}
\newcommand{\Tsolver}{\ensuremath{\mathsf{SOLVER}}}
\newcommand{\Tproof}{\ensuremath{\mathsf{PROOF}}}

% Site and geometry
\newcommand{\Site}{\mathcal{S}}             % semantic site
\newcommand{\Cov}{\mathrm{Cov}}            % covering family
\newcommand{\Ob}{\mathrm{Ob}}              % objects
\newcommand{\Mor}{\mathrm{Mor}}            % morphisms
\newcommand{\res}{\mathrm{res}}            % restriction
\newcommand{\Cech}{\check{C}}              % Čech complex
\newcommand{\Hcoh}[1]{H^{#1}}             % cohomology group

% Descent
\newcommand{\descent}{\mathrm{desc}}
\newcommand{\glue}{\mathrm{glue}}
\newcommand{\overlap}[2]{#1 \cap #2}

% Semantic moves
\newcommand{\VERIFY}{\textsc{Verify}}
\newcommand{\CONSTRUCT}{\textsc{Construct}}
\newcommand{\NORMALIZE}{\textsc{Normalize}}
\newcommand{\GLUE}{\textsc{Glue}}
\newcommand{\RESTRICT}{\textsc{Restrict}}
\newcommand{\TRANSPORT}{\textsc{Transport}}
\newcommand{\REPAIR}{\textsc{Repair}}
\newcommand{\PROMOTE}{\textsc{Promote}}
\newcommand{\CHALLENGE}{\textsc{Challenge}}
\newcommand{\ESCALATE}{\textsc{Escalate}}

% Fragments
\newcommand{\QFLIA}{\texttt{QF\_LIA}}
\newcommand{\QFLRA}{\texttt{QF\_LRA}}
\newcommand{\QFBV}{\texttt{QF\_BV}}
\newcommand{\QFUF}{\texttt{QF\_UF}}

% Misc
\newcommand{\Python}{\textsc{Python}}
\newcommand{\JuGeo}{\textsc{JuGeo}}
\newcommand{\LEAN}{\textsc{Lean}\,4}
\newcommand{\Fstar}{F$^{\star}$}
\newcommand{\Coq}{\textsc{Coq}}
\newcommand{\Dafny}{\textsc{Dafny}}

% Common shortcuts
\newcommand{\ie}{i.e.\@\xspace}
\newcommand{\eg}{e.g.\@\xspace}
\newcommand{\cf}{cf.\@\xspace}
\newcommand{\wrt}{w.r.t.\@\xspace}

% Evaluation shortcuts
\newcommand{\numcases}{300}
\newcommand{\numspec}{100}
\newcommand{\numeq}{100}
\newcommand{\numbug}{100}
\newcommand{\medtime}{6.5\,ms}
\newcommand{\totaltime}{1.949\,s}


% ── Packages ────────────────────────────────────────────────────────────
\usepackage{amsmath,amssymb,amsthm,mathtools}
\usepackage{tikz-cd}
\usepackage{listings}
\usepackage{xcolor}
\usepackage{xspace}
\usepackage{booktabs}
\usepackage{multirow}
\usepackage{enumitem}
\usepackage[expansion=false]{microtype}
\usepackage{hyperref}
\usepackage{cleveref}
% \usepackage{thmtools}  % disabled: not available in this TeX installation
\usepackage{float}
\usepackage{caption}
\usepackage{subcaption}
\usepackage{tabularx}
\usepackage{stmaryrd}       % \llbracket, \rrbracket
\usepackage{bbm}            % \mathbbm{1}

% ── Page geometry ───────────────────────────────────────────────────────
\usepackage[margin=1in]{geometry}

% ── Theorem environments ────────────────────────────────────────────────
\theoremstyle{plain}
\newtheorem{theorem}{Theorem}[section]
\newtheorem{lemma}[theorem]{Lemma}
\newtheorem{proposition}[theorem]{Proposition}
\newtheorem{corollary}[theorem]{Corollary}

\theoremstyle{definition}
\newtheorem{definition}[theorem]{Definition}
\newtheorem{example}[theorem]{Example}
\newtheorem{construction}[theorem]{Construction}

\theoremstyle{remark}
\newtheorem{remark}[theorem]{Remark}
\newtheorem{notation}[theorem]{Notation}

% ── Python listings style ──────────────────────────────────────────────
\definecolor{jugeoBlue}{HTML}{2563EB}
\definecolor{jugeoGreen}{HTML}{059669}
\definecolor{jugeoRed}{HTML}{DC2626}
\definecolor{jugeoPurple}{HTML}{7C3AED}
\definecolor{jugeoGray}{HTML}{6B7280}
\definecolor{jugeoBg}{HTML}{F8FAFC}

\lstdefinestyle{jugeo-python}{
  language=Python,
  basicstyle=\ttfamily\small,
  keywordstyle=\color{jugeoBlue}\bfseries,
  stringstyle=\color{jugeoGreen},
  commentstyle=\color{jugeoGray}\itshape,
  numberstyle=\tiny\color{jugeoGray},
  backgroundcolor=\color{jugeoBg},
  numbers=left,
  numbersep=6pt,
  frame=single,
  framerule=0.4pt,
  rulecolor=\color{jugeoGray!40},
  breaklines=true,
  breakatwhitespace=true,
  showstringspaces=false,
  tabsize=4,
  xleftmargin=1.5em,
  framexleftmargin=1.5em,
  morekeywords={dataclass,frozen,field,Enum,IntEnum,auto,Any,
                Mapping,Sequence,Optional,match,case},
  literate={→}{{$\to$}}1 {←}{{$\leftarrow$}}1
           {≤}{{$\leq$}}1 {≥}{{$\geq$}}1 {≠}{{$\neq$}}1
           {∈}{{$\in$}}1 {∀}{{$\forall$}}1 {∃}{{$\exists$}}1
           {⊕}{{$\oplus$}}1 {⊖}{{$\ominus$}}1,
  escapeinside={(*@}{@*)},
}
\lstset{style=jugeo-python}

\lstdefinestyle{jugeo-output}{
  basicstyle=\ttfamily\small,
  backgroundcolor=\color{jugeoBg},
  frame=single,
  framerule=0.4pt,
  rulecolor=\color{jugeoGray!40},
  breaklines=true,
  showstringspaces=false,
  xleftmargin=1.5em,
  framexleftmargin=0.5em,
  numbers=none,
}

% ── JuGeo-specific macros ──────────────────────────────────────────────

% Judgment 8-tuple
\newcommand{\J}{\mathcal{J}}
\newcommand{\Jud}[8]{(#1,\,#2,\,#3,\,#4,\,#5,\,#6,\,#7,\,#8)}
\newcommand{\JudShort}[1]{\J_{#1}}

% Components
\newcommand{\coord}{\mathsf{c}}            % coordinate
\newcommand{\prop}{\varphi}                 % proposition
\newcommand{\carrier}{\mathcal{A}}          % carrier type
\newcommand{\evid}{\mathcal{E}}             % evidence bundle
\newcommand{\oblig}{\mathcal{O}}            % obligations
\newcommand{\obstr}{\mathcal{B}}            % obstructions
\newcommand{\trust}{\mathcal{T}}            % trust annotation
\newcommand{\prov}{\Pi}                     % provenance

% Trust algebra
\newcommand{\Talg}{\mathfrak{T}}            % trust ordered algebra
\newcommand{\Eadm}{\mathcal{E}_{\mathrm{adm}}}
\newcommand{\tleq}{\preceq}                 % trust ordering
\newcommand{\tjoin}{\oplus}                 % conservative join
\newcommand{\tatten}{\ominus}               % attenuation
\newcommand{\tpromote}{\uparrow_\pi}        % promotion
\newcommand{\tdemote}{\downarrow_\chi}       % demotion

% Trust tiers
\newcommand{\Tcontra}{\ensuremath{\mathsf{CONTRADICTED}}}
\newcommand{\Tunver}{\ensuremath{\mathsf{UNVERIFIED}}}
\newcommand{\Tcopilot}{\ensuremath{\mathsf{COPILOT}}}
\newcommand{\Toracle}{\ensuremath{\mathsf{ORACLE}}}
\newcommand{\Truntime}{\ensuremath{\mathsf{RUNTIME}}}
\newcommand{\Tsolver}{\ensuremath{\mathsf{SOLVER}}}
\newcommand{\Tproof}{\ensuremath{\mathsf{PROOF}}}

% Site and geometry
\newcommand{\Site}{\mathcal{S}}             % semantic site
\newcommand{\Cov}{\mathrm{Cov}}            % covering family
\newcommand{\Ob}{\mathrm{Ob}}              % objects
\newcommand{\Mor}{\mathrm{Mor}}            % morphisms
\newcommand{\res}{\mathrm{res}}            % restriction
\newcommand{\Cech}{\check{C}}              % Čech complex
\newcommand{\Hcoh}[1]{H^{#1}}             % cohomology group

% Descent
\newcommand{\descent}{\mathrm{desc}}
\newcommand{\glue}{\mathrm{glue}}
\newcommand{\overlap}[2]{#1 \cap #2}

% Semantic moves
\newcommand{\VERIFY}{\textsc{Verify}}
\newcommand{\CONSTRUCT}{\textsc{Construct}}
\newcommand{\NORMALIZE}{\textsc{Normalize}}
\newcommand{\GLUE}{\textsc{Glue}}
\newcommand{\RESTRICT}{\textsc{Restrict}}
\newcommand{\TRANSPORT}{\textsc{Transport}}
\newcommand{\REPAIR}{\textsc{Repair}}
\newcommand{\PROMOTE}{\textsc{Promote}}
\newcommand{\CHALLENGE}{\textsc{Challenge}}
\newcommand{\ESCALATE}{\textsc{Escalate}}

% Fragments
\newcommand{\QFLIA}{\texttt{QF\_LIA}}
\newcommand{\QFLRA}{\texttt{QF\_LRA}}
\newcommand{\QFBV}{\texttt{QF\_BV}}
\newcommand{\QFUF}{\texttt{QF\_UF}}

% Misc
\newcommand{\Python}{\textsc{Python}}
\newcommand{\JuGeo}{\textsc{JuGeo}}
\newcommand{\LEAN}{\textsc{Lean}\,4}
\newcommand{\Fstar}{F$^{\star}$}
\newcommand{\Coq}{\textsc{Coq}}
\newcommand{\Dafny}{\textsc{Dafny}}

% Common shortcuts
\newcommand{\ie}{i.e.\@\xspace}
\newcommand{\eg}{e.g.\@\xspace}
\newcommand{\cf}{cf.\@\xspace}
\newcommand{\wrt}{w.r.t.\@\xspace}

% Evaluation shortcuts
\newcommand{\numcases}{300}
\newcommand{\numspec}{100}
\newcommand{\numeq}{100}
\newcommand{\numbug}{100}
\newcommand{\medtime}{6.5\,ms}
\newcommand{\totaltime}{1.949\,s}

% No external data file for this paper.

% ── Additional packages ────────────────────────────────────────────
\usepackage{array}
\usepackage{tikz}
\usetikzlibrary{arrows.meta,positioning,calc,fit,backgrounds}

% ── Paper-specific macros ──────────────────────────────────────────
\newcommand{\CopChan}{\mathsf{CopilotChannel}}
\newcommand{\EvChan}{\mathsf{EvidenceChannel}}
\newcommand{\ChanConfig}{\mathsf{ChannelConfiguration}}
\newcommand{\ChanJuris}{\mathsf{ChannelJurisdiction}}
\newcommand{\ChanPool}{\mathsf{ChannelPool}}
\newcommand{\ChanMon}{\mathsf{ChannelMonitor}}
\newcommand{\ChanSer}{\mathsf{ChannelSerializer}}
\newcommand{\ChanRouter}{\mathsf{ChannelRouter}}
\newcommand{\ChanFed}{\mathsf{ChannelFederation}}
\newcommand{\EvReq}{\mathsf{EvidenceRequest}}
\newcommand{\EvResp}{\mathsf{EvidenceResponse}}
\newcommand{\EvRec}{\mathsf{EvidenceRecord}}
\newcommand{\TrustCeil}{\mathrm{ceiling}}
\newcommand{\CopQuery}{\texttt{query\_llm}}
\newcommand{\CopParse}{\texttt{parse\_response}}
\newcommand{\CopHealth}{\texttt{health\_check}}
\newcommand{\CopHandle}{\texttt{handle\_request}}
\newcommand{\EffCeil}{\texttt{effective\_trust\_ceiling}}
\newcommand{\AdmCheck}{\texttt{check\_admissibility}}
\newcommand{\RoundTrip}{\texttt{round\_trip}}

\title{\textbf{The Copilot Channel:\\
Architecture, Jurisdiction, and Infrastructure}}

\author{Halley Young}

\date{Paper~76 --- Judgment Geometry Series}

\begin{document}
\maketitle

% ─────────────────────────────────────────────────────────────────────
\begin{abstract}
The $\CopChan$ is the primary conduit through which LLM-generated
evidence enters the \JuGeo{} verification pipeline.  Unlike solver
or runtime channels, the Copilot channel produces evidence with no
intrinsic soundness guarantee, necessitating a carefully designed
trust ceiling, jurisdiction policy, and infrastructure layer.  In
this paper we present a comprehensive study of the $\CopChan$ class
architecture, covering:
(i)~the $\EvChan$ enumeration and the channel taxonomy that positions
Copilot as the unique \emph{mechanical-yet-corroboration-required}
channel;
(ii)~the $\ChanJuris$ mechanism that confines LLM evidence to
specified domains and proposition kinds;
(iii)~the $\ChanConfig$ dataclass and its interaction with rate
limiting, timeouts, and trust ceilings;
(iv)~the $\ChanPool$, $\ChanMon$, and $\ChanSer$ infrastructure
classes that manage channel lifecycle, performance monitoring, and
serialization.

We prove five formal properties: jurisdiction soundness, trust-ceiling
idempotence, rate-limit correctness, serialization round-trip
fidelity, and monitor completeness.  All results are proved in full
and formalized in \LEAN.  We reference
\pSevenSixQueries{} verification queries spanning
\pSevenSixUniqueRepos{} repositories and
\pSevenSixLanguages{} programming languages.
\end{abstract}

% ─────────────────────────────────────────────────────────────────────
\section{Introduction}
\label{sec:intro}

In the \JuGeo{} framework, evidence channels are the structured
conduits through which different sources of verification evidence
contribute to the judgment state.  Each channel is characterized by
its trust floor, mechanicality, corroboration requirements, and
jurisdiction.  The $\CopChan$ occupies a unique position in this
taxonomy: it is the only channel that is both \emph{mechanical}
(fully automated, no human in the loop) and \emph{requires
corroboration} (its evidence alone cannot elevate trust above
$\Tcopilot$).

This dual nature reflects the fundamental characteristic of
LLM-generated evidence: it is produced algorithmically and at
scale, but without the formal guarantees provided by SMT solvers
or proof checkers.  Integrating such evidence into a
trust-calibrated verification system requires:

\begin{itemize}[leftmargin=2em]
\item A \emph{trust ceiling} that prevents LLM evidence from
  claiming higher trust than warranted.
\item A \emph{jurisdiction policy} that confines LLM evidence to
  appropriate domains and proposition kinds.
\item A \emph{rate-limiting mechanism} to prevent resource exhaustion.
\item A \emph{monitoring infrastructure} to track channel performance
  and detect degradation.
\item A \emph{serialization layer} for persistent storage and
  reproducibility.
\end{itemize}

\paragraph{Contributions.}
The contributions of this paper are:
\begin{itemize}[leftmargin=2em]
\item We present the full architecture of the $\CopChan$ class,
  including its configuration, jurisdiction, and trust enforcement
  mechanisms (\S\ref{sec:arch}).

\item We formalize the jurisdiction mechanism and prove that it
  correctly confines evidence to admissible domains
  (\S\ref{sec:jurisdiction}).

\item We describe the infrastructure layer ($\ChanPool$, $\ChanMon$,
  $\ChanSer$) and prove key properties of each component
  (\S\ref{sec:infra}).

\item We prove five formal properties: jurisdiction soundness,
  ceiling idempotence, rate-limit correctness, round-trip fidelity,
  and monitor completeness (\S\ref{sec:formal}).
\end{itemize}

\paragraph{Outline.}
Section~\ref{sec:arch} describes the channel architecture.
Section~\ref{sec:impl} details the $\CopChan$ implementation.
Section~\ref{sec:jurisdiction} formalizes jurisdiction and configuration.
Section~\ref{sec:infra} covers the infrastructure layer.
Section~\ref{sec:formal} states and proves the main theorems.
Section~\ref{sec:discussion} discusses design decisions.
Section~\ref{sec:related} surveys related work.
Section~\ref{sec:conclusion} concludes.

% ─────────────────────────────────────────────────────────────────────
\section{Channel Architecture}
\label{sec:arch}

\subsection{The EvidenceChannel Enumeration}
\label{sec:arch:enum}

The $\EvChan$ enumeration defines the seven evidence channels
supported by \JuGeo{}:

\begin{definition}[Evidence channel enumeration]
\label{def:ev-channel}
The $\EvChan$ enumeration consists of:
$\{\mathsf{SOLVER}, \mathsf{RUNTIME}, \mathsf{ORACLE},
\mathsf{COPILOT}, \mathsf{FORMAL\_PROOF}, \mathsf{HUMAN},
\mathsf{COMPOSED}\}$.
Each member $C$ has:
\begin{itemize}[leftmargin=2em]
\item A \emph{trust floor} $\tau_{\min}(C)$: the minimum trust
  this channel may assign.
\item A \emph{mechanicality flag} $\mathrm{mech}(C) \in \{\top, \bot\}$.
\item A \emph{corroboration flag} $\mathrm{corrob}(C) \in \{\top, \bot\}$.
\item A set of \emph{default query families} $\mathrm{QF}(C)$.
\end{itemize}
\end{definition}

\begin{remark}[Copilot's unique position]
\label{rem:copilot-unique}
The $\mathsf{COPILOT}$ channel is the only channel with
$\mathrm{mech}(\mathsf{COPILOT}) = \top$ and
$\mathrm{corrob}(\mathsf{COPILOT}) = \top$.  This means it
operates without human involvement but its evidence requires
independent confirmation.  The $\mathsf{SOLVER}$ and
$\mathsf{FORMAL\_PROOF}$ channels are mechanical but do not
require corroboration (their evidence is sound by construction).
The $\mathsf{ORACLE}$ and $\mathsf{HUMAN}$ channels require
corroboration but are not mechanical.
\end{remark}

\subsection{Channel Taxonomy}
\label{sec:arch:taxonomy}

The seven channels partition into four groups based on
mechanicality and corroboration:

\begin{definition}[Channel classification]
\label{def:channel-class}
\begin{enumerate}
\item \textbf{Mechanical, no corroboration:} $\{\mathsf{SOLVER},
  \mathsf{RUNTIME}, \mathsf{FORMAL\_PROOF}\}$.  Sound by
  construction; evidence can be directly trusted.
\item \textbf{Mechanical, requires corroboration:}
  $\{\mathsf{COPILOT}\}$.  Automated but not sound; evidence
  needs independent confirmation.
\item \textbf{Non-mechanical, requires corroboration:}
  $\{\mathsf{ORACLE}, \mathsf{HUMAN}\}$.  Human-in-the-loop;
  evidence needs confirmation.
\item \textbf{Non-mechanical, no corroboration:}
  $\emptyset$.  This category is intentionally empty in \JuGeo{};
  non-mechanical evidence always requires corroboration.
\end{enumerate}
\end{definition}

\begin{lemma}[Corroboration monotonicity]
\label{lem:corrob-mono}
If $\mathrm{corrob}(C) = \bot$ then $\tau_{\min}(C) \tleq \Tsolver$
(\ie only high-trust channels are exempt from corroboration).
\end{lemma}

\begin{proof}
By inspection of the channel taxonomy: the channels with
$\mathrm{corrob} = \bot$ are $\mathsf{SOLVER}$,
$\mathsf{RUNTIME}$, and $\mathsf{FORMAL\_PROOF}$, with trust
floors $\Tsolver$, $\Truntime$, and $\Tproof$ respectively.
All satisfy $\tau_{\min}(C) \tleq \Tsolver$.
\end{proof}

\subsection{Trust Floors and the Trust Lattice}
\label{sec:arch:trust}

\begin{definition}[Channel trust interval]
\label{def:trust-interval}
For channel $C$ with trust floor $\tau_{\min}(C)$ and
ceiling $\tau_{\max}(C)$, the \emph{trust interval} is
$I(C) = [\tau_{\min}(C), \tau_{\max}(C)]$.
For the $\CopChan$: $I(\mathsf{COPILOT}) = [\Tcopilot, \Tcopilot]$.
\end{definition}

The trust interval of the Copilot channel is a singleton:
evidence is always assigned exactly $\Tcopilot$ (unless
corroboration promotes it, which is handled outside the channel).

The following listing shows the complete setup of a $\CopChan$
with its configuration.

\begin{lstlisting}[style=jugeo-python, caption={Complete CopilotChannel configuration and setup.}, label=lst:config]
from jugeo.evidence.channels import (
    CopilotChannel,
    ChannelConfiguration,
    EvidenceChannel,
    ChannelJurisdiction,
    EvidenceRequest,
    EvidenceResponse,
    EvidenceKind,
)

# Build jurisdiction for the Copilot channel
jurisdiction = ChannelJurisdiction.for_channel(
    EvidenceChannel.COPILOT
)

# Full configuration with all parameters
config = ChannelConfiguration(
    channel=EvidenceChannel.COPILOT,
    timeout_ms=10000,
    max_retries=3,
    batch_size=1,
    rate_limit=10.0,
    trust_ceiling="proposal",
    jurisdiction=jurisdiction,
    is_enabled=True,
    priority=40,
)

# Verify configuration properties
assert config.channel == EvidenceChannel.COPILOT
eff_ceil = config.effective_trust_ceiling()
print(f"Effective trust ceiling: {eff_ceil}")

# Create a channel with modified timeout
fast_config = config.with_timeout(5000)
assert fast_config.timeout_ms == 5000

# Default configuration shorthand
default_config = ChannelConfiguration.default_for(
    EvidenceChannel.COPILOT
)

# Instantiate the channel
copilot = CopilotChannel(
    config=config,
    model_name="copilot-default",
)

# Verify the trust ceiling constant
assert copilot.TRUST_CEILING == "proposal"

# Check channel health
health = copilot.health_check()
print(f"Channel healthy: {health}")

# Query the LLM
result = copilot.query_llm(
    prompt="Suggest an invariant for: while i < n: sum += a[i]; i += 1",
    timeout_ms=8000,
)

# Parse and apply trust ceiling
parsed = copilot.parse_response()
copilot.apply_trust_ceiling()

# Handle a full evidence request
request = EvidenceRequest(
    request_id="chan-test-001",
    coordinate="loop.sum.invariant",
    proposition="sum == partial_sum(a, i)",
    required_kind="PROPOSAL",
    proposition_kind="invariant",
    preferred_channel=EvidenceChannel.COPILOT,
    fallback_channels=(EvidenceChannel.SOLVER,),
    deadline_ms=10000.0,
    budget=1.0,
    metadata={"source": "paper76-example"},
)

response = copilot.handle_request(request)
assert response.trust_level == "proposal"
print(f"Response partial: {response.is_partial}")
print(f"Response empty: {response.is_empty()}")
clamped = response.clamp_trust("proposal")
print(f"Clamped trust: {clamped.trust_level}")
\end{lstlisting}

% ─────────────────────────────────────────────────────────────────────
\section{CopilotChannel Implementation}
\label{sec:impl}

The $\CopChan$ class implements the channel interface with
specializations for LLM-generated evidence.

\subsection{Class Structure}
\label{sec:impl:structure}

The $\CopChan$ has the following key attributes:
\begin{itemize}[leftmargin=2em]
\item $\texttt{TRUST\_CEILING} = \texttt{"proposal"}$: a class-level
  constant defining the maximum trust.
\item $\texttt{config}$: a $\ChanConfig$ instance controlling
  timeout, retries, rate limit, and jurisdiction.
\item $\texttt{model\_name}$: the identifier of the underlying LLM.
\end{itemize}

\begin{definition}[CopilotChannel specification]
\label{def:copilot-spec}
The $\CopChan$ is specified by the tuple
$(C, M, \tau_c, J, \lambda, \rho)$ where:
\begin{itemize}[leftmargin=2em]
\item $C = \mathsf{COPILOT}$ is the channel identifier.
\item $M$ is the model name string.
\item $\tau_c = \Tcopilot$ is the trust ceiling.
\item $J$ is the jurisdiction (a $\ChanJuris$ instance).
\item $\lambda$ is the rate limit (queries per second).
\item $\rho$ is the maximum retry count.
\end{itemize}
\end{definition}

\subsection{The query\_llm Method}
\label{sec:impl:query}

The \CopQuery{} method is the primary interface for LLM inference:

\begin{definition}[LLM query specification]
\label{def:query-spec}
The $\CopQuery(\pi, \delta)$ method takes a prompt $\pi$ and
optional timeout $\delta$ and returns a dictionary $R$ with:
\begin{itemize}[leftmargin=2em]
\item $R[\texttt{text}]$: the raw LLM response text.
\item $R[\texttt{model}]$: the model identifier.
\item $R[\texttt{latency\_ms}]$: the inference latency.
\item $R[\texttt{tokens}]$: the token count.
\item $R[\texttt{truncated}]$: whether the response was truncated.
\end{itemize}
The method respects the timeout bound: if the LLM does not respond
within $\delta$ milliseconds, the query is aborted and an empty
result is returned.
\end{definition}

\subsection{Trust Ceiling Enforcement}
\label{sec:impl:ceiling}

Trust ceiling enforcement is a two-stage process:

\begin{enumerate}
\item \textbf{Channel-level enforcement.}  The
  $\texttt{apply\_trust\_ceiling}()$ method clamps the trust of
  the parsed response to at most $\texttt{TRUST\_CEILING}$.

\item \textbf{Bridge-level enforcement.}  The $\CopBridge$
  (described in Paper~93) applies a second enforcement layer
  via $\texttt{enforce\_trust\_ceiling}()$.
\end{enumerate}

\begin{definition}[Ceiling enforcement function]
\label{def:enforce}
The ceiling enforcement function
$\TrustCeil_c : \mathcal{T} \to \mathcal{T}$ is:
\[
  \TrustCeil_c(\tau) = \begin{cases}
    \tau & \text{if } \tau \tleq c, \\
    c & \text{otherwise.}
  \end{cases}
\]
For the $\CopChan$, $c = \Tcopilot$.
\end{definition}

\begin{lemma}[Ceiling enforcement is idempotent]
\label{lem:ceil-idemp}
$\TrustCeil_c(\TrustCeil_c(\tau)) = \TrustCeil_c(\tau)$ for all
$\tau, c$.
\end{lemma}

\begin{proof}
Let $\tau' = \TrustCeil_c(\tau)$.  By definition, $\tau' \tleq c$.
Therefore $\TrustCeil_c(\tau') = \tau'$ (the first case applies).
\end{proof}

% ─────────────────────────────────────────────────────────────────────
\section{Jurisdiction and Configuration}
\label{sec:jurisdiction}

\subsection{ChannelJurisdiction}
\label{sec:jurisdiction:juris}

The $\ChanJuris$ controls which evidence requests the $\CopChan$
may handle.

\begin{definition}[Channel jurisdiction]
\label{def:jurisdiction}
A \emph{channel jurisdiction} $J$ is a tuple
$(D, P, K, \tau_{\max}, r, X)$ where:
\begin{itemize}[leftmargin=2em]
\item $D$ is the set of admissible domain strings.
\item $P$ is the set of coordinate patterns (glob-like patterns).
\item $K$ is the set of admissible proposition kinds.
\item $\tau_{\max}$ is the maximum trust level for this jurisdiction.
\item $r \in \{\top, \bot\}$ is the corroboration requirement.
\item $X$ is the set of excluded query families.
\end{itemize}
\end{definition}

\begin{definition}[Admissibility]
\label{def:admissibility}
A request $q = (c, \varphi, k)$ is \emph{admissible} under
jurisdiction $J = (D, P, K, \tau_{\max}, r, X)$ if:
\begin{enumerate}
\item The coordinate $c$ matches at least one pattern in $P$.
\item The proposition kind $k$ is in $K$.
\item The query family of $q$ is not in $X$.
\end{enumerate}
We write $J \vdash q$ for admissibility and $J \nvdash q$ for
non-admissibility.
\end{definition}

The following listing demonstrates jurisdiction setup and
admissibility checking.

\begin{lstlisting}[style=jugeo-python, caption={Jurisdiction configuration and admissibility checking.}, label=lst:jurisdiction]
from jugeo.evidence.channels import (
    ChannelJurisdiction,
    ChannelConfiguration,
    EvidenceChannel,
    EvidenceRequest,
    EvidenceKind,
)

# Create jurisdiction for the Copilot channel
copilot_juris = ChannelJurisdiction.for_channel(
    EvidenceChannel.COPILOT
)

# Inspect jurisdiction properties
print(f"Domain set: {copilot_juris.domain_set}")
print(f"Coordinate patterns: {copilot_juris.coordinate_patterns}")
print(f"Proposition kinds: {copilot_juris.proposition_kinds}")
print(f"Max trust level: {copilot_juris.max_trust_level}")
print(f"Requires corroboration: {copilot_juris.requires_corroboration}")
print(f"Excluded families: {copilot_juris.excluded_families}")

# Check admissibility of a request
admissible = copilot_juris.check_admissibility(
    coordinate="module.function.invariant",
    proposition_kind="invariant",
)
print(f"Admissible: {admissible}")

# Check all conditions at once
result = copilot_juris.check_all()
print(f"Jurisdiction check: {result}")

# Create jurisdiction for a different channel for comparison
solver_juris = ChannelJurisdiction.for_channel(
    EvidenceChannel.SOLVER
)

# Create jurisdiction by kind
proposal_juris = ChannelJurisdiction.for_kind(
    EvidenceKind.PROPOSAL
)

# Serialize jurisdiction
juris_dict = copilot_juris.to_dict()
restored = ChannelJurisdiction.from_dict(juris_dict)

# Build full configuration with jurisdiction
config = ChannelConfiguration(
    channel=EvidenceChannel.COPILOT,
    timeout_ms=8000,
    max_retries=2,
    batch_size=1,
    rate_limit=5.0,
    trust_ceiling="proposal",
    jurisdiction=copilot_juris,
    is_enabled=True,
    priority=45,
)

# Configuration serialization
config_dict = config.to_dict()
print(f"Config serialized: {list(config_dict.keys())}")
\end{lstlisting}

\subsection{ChannelConfiguration}
\label{sec:jurisdiction:config}

\begin{definition}[Channel configuration]
\label{def:config}
A \emph{channel configuration} $\kappa$ is a dataclass
$(\EvChan, \delta, \rho, \beta, \lambda, \tau_c, J, e, p)$ where:
\begin{itemize}[leftmargin=2em]
\item $\EvChan$ is the channel identifier.
\item $\delta$ is the timeout in milliseconds.
\item $\rho$ is the maximum retry count.
\item $\beta$ is the batch size.
\item $\lambda$ is the rate limit (requests/second, 0 = unlimited).
\item $\tau_c$ is the trust ceiling string.
\item $J$ is the jurisdiction.
\item $e$ is the enabled flag.
\item $p$ is the priority (lower = higher priority).
\end{itemize}
\end{definition}

\begin{definition}[Effective trust ceiling]
\label{def:eff-ceil}
The \emph{effective trust ceiling} of a configuration $\kappa$ is:
\[
  \EffCeil(\kappa) = \min(\tau_c(\kappa), \tau_{\max}(J(\kappa)))
\]
\ie the minimum of the configured ceiling and the jurisdiction's
maximum trust.  This ensures that the jurisdiction cannot weaken
the ceiling but may strengthen it.
\end{definition}

\begin{lemma}[Effective ceiling is conservative]
\label{lem:eff-ceil-conservative}
$\EffCeil(\kappa) \tleq \tau_c(\kappa)$ and
$\EffCeil(\kappa) \tleq \tau_{\max}(J(\kappa))$.
\end{lemma}

\begin{proof}
Immediate from the definition of $\min$.
\end{proof}

% ─────────────────────────────────────────────────────────────────────
\section{Channel Infrastructure}
\label{sec:infra}

The $\CopChan$ operates within a broader infrastructure that
manages channel lifecycle, monitors performance, and handles
serialization.

\subsection{ChannelPool}
\label{sec:infra:pool}

The $\ChanPool$ manages the set of active channels.

\begin{definition}[Channel pool]
\label{def:pool}
A \emph{channel pool} $\Pi$ is a mapping from channel identifiers
to channel instances: $\Pi : \EvChan \to \mathrm{Channel}$.
The pool supports:
\begin{itemize}[leftmargin=2em]
\item $\texttt{register\_channel}(c)$: add channel $c$ to the pool.
\item $\texttt{get\_channel}(id)$: retrieve channel by identifier.
\item $\texttt{has\_channel}(id)$: check presence.
\item $\texttt{health\_check\_all}()$: check health of all channels.
\item $\texttt{drain}()$: gracefully shut down all channels.
\end{itemize}
\end{definition}

\subsection{ChannelMonitor}
\label{sec:infra:monitor}

The $\ChanMon$ tracks channel performance metrics.

\begin{definition}[Channel monitor]
\label{def:monitor}
A \emph{channel monitor} $\mu$ maintains:
\begin{itemize}[leftmargin=2em]
\item A request log $L_R$: records of all requests.
\item A response log $L_S$: records of all responses.
\item A failure log $L_F$: records of failures.
\item Derived metrics: latency percentiles, success rate, throughput.
\end{itemize}
\end{definition}

\begin{definition}[Monitor completeness]
\label{def:monitor-complete}
A monitor $\mu$ is \emph{complete} for a sequence of channel
interactions $I_1, \ldots, I_n$ if for every interaction $I_k$:
\begin{enumerate}
\item If $I_k$ is a request, there exists a record in $L_R$.
\item If $I_k$ is a response, there exists a record in $L_S$.
\item If $I_k$ is a failure, there exists a record in $L_F$.
\end{enumerate}
\end{definition}

\subsection{ChannelSerializer}
\label{sec:infra:serializer}

The $\ChanSer$ provides serialization and deserialization for
channel messages.

\begin{definition}[Serialization round-trip]
\label{def:round-trip}
A serializer $S$ has \emph{round-trip fidelity} if for every
request $r$ and response $s$:
\[
  S.\texttt{from\_json}(S.\texttt{to\_json}(r)) = r
  \quad\text{and}\quad
  S.\texttt{from\_json}(S.\texttt{to\_json}(s)) = s.
\]
\end{definition}

The following listing demonstrates the channel infrastructure
in action.

\begin{lstlisting}[style=jugeo-python, caption={Channel pool, monitor, and serializer infrastructure.}, label=lst:infra]
from jugeo.evidence.channels import (
    CopilotChannel,
    SolverChannel,
    RuntimeChannel,
    ChannelPool,
    ChannelMonitor,
    ChannelSerializer,
    ChannelConfiguration,
    EvidenceChannel,
    EvidenceRequest,
    EvidenceResponse,
)

# Set up channel pool with multiple channels
pool = ChannelPool()

copilot_config = ChannelConfiguration.default_for(
    EvidenceChannel.COPILOT
)
solver_config = ChannelConfiguration.default_for(
    EvidenceChannel.SOLVER
)

copilot = CopilotChannel(config=copilot_config)
solver = SolverChannel(config=solver_config)
runtime = RuntimeChannel()

pool.register_channel(copilot)
pool.register_channel(solver)
pool.register_channel(runtime)

# Check pool state
assert pool.has_channel(EvidenceChannel.COPILOT)
assert pool.has_channel(EvidenceChannel.SOLVER)

# Health check all channels
health_results = pool.health_check_all()
print(f"Pool health: {health_results}")

# Set up monitoring
monitor = ChannelMonitor()

# Process a request through the pool
request = EvidenceRequest(
    request_id="pool-test-001",
    coordinate="module.function.spec",
    proposition="result >= 0",
    required_kind="PROPOSAL",
    proposition_kind="postcondition",
    preferred_channel=EvidenceChannel.COPILOT,
    fallback_channels=(EvidenceChannel.SOLVER,),
    deadline_ms=5000.0,
    budget=1.0,
    metadata={},
)

# Record and route the request
monitor.record_request(request)
channel = pool.get_channel(EvidenceChannel.COPILOT)
response = channel.handle_request(request)
monitor.record_response(response)

# Query monitor metrics
print(f"Success rate: {monitor.success_rate()}")
print(f"Throughput: {monitor.throughput()}")
percentiles = monitor.latency_percentiles()
print(f"Latency percentiles: {percentiles}")
summary = monitor.summary()
print(f"Monitor summary: {summary}")

# Serialization round-trip
req_json = ChannelSerializer.request_to_json(request)
resp_json = ChannelSerializer.response_to_json(response)

# Verify round-trip fidelity
rt_request = ChannelSerializer.round_trip_request(request)
rt_response = ChannelSerializer.round_trip_response(response)
print(f"Request round-trip OK: {rt_request.request_id == request.request_id}")
print(f"Response round-trip OK: {rt_response.request_id == response.request_id}")

# Graceful drain
if pool.is_draining():
    pool.drain()
\end{lstlisting}

\subsection{Channel Router Integration}
\label{sec:infra:router}

The $\ChanRouter$ determines which channel handles each request.

\begin{definition}[Routing policy]
\label{def:routing}
The routing policy $R : \EvReq \to \EvChan$ maps each request to
a channel based on:
\begin{enumerate}
\item The preferred channel (if specified and available).
\item The jurisdiction of each channel (only admissible channels
  are candidates).
\item The priority ordering of candidate channels.
\item The fallback chain (if the preferred channel is unavailable
  or inadmissible).
\end{enumerate}
\end{definition}

The following listing shows the complete serialization round-trip
for evidence requests and responses.

\begin{lstlisting}[style=jugeo-python, caption={Complete serialization round-trip demonstration.}, label=lst:serializer]
from jugeo.evidence.channels import (
    ChannelSerializer,
    EvidenceRequest,
    EvidenceResponse,
    EvidenceChannel,
    ChannelJurisdiction,
    ChannelConfiguration,
)

# Create a request to serialize
request = EvidenceRequest(
    request_id="serialize-test-001",
    coordinate="module.class.method.precondition",
    proposition="x > 0 and y > 0 implies x + y > 0",
    required_kind="PROPOSAL",
    proposition_kind="precondition",
    preferred_channel=EvidenceChannel.COPILOT,
    fallback_channels=(
        EvidenceChannel.SOLVER,
        EvidenceChannel.RUNTIME,
    ),
    deadline_ms=15000.0,
    budget=2.0,
    metadata={
        "source": "paper76",
        "priority": "high",
    },
)

# Serialize to JSON
request_json = ChannelSerializer.request_to_json(request)
print(f"Request JSON keys: {list(request_json.keys())}")

# Deserialize back
restored_request = ChannelSerializer.request_from_json(request_json)
assert restored_request.request_id == request.request_id
assert restored_request.coordinate == request.coordinate
assert restored_request.proposition == request.proposition

# Canonical key for deduplication
key = request.canonical_key()
print(f"Canonical key: {key}")

# Create a response to serialize
response = EvidenceResponse(
    request_id="serialize-test-001",
    response_id="resp-001",
    channel=EvidenceChannel.COPILOT,
    evidence_item={"type": "invariant", "text": "x > 0"},
    evidence={"raw": "suggested invariant: x > 0"},
    trust_level="proposal",
    latency_ms=245.3,
    is_partial=False,
    residuals=("verify_boundary",),
    provenance=("copilot-default",),
)

# Serialize response
response_json = ChannelSerializer.response_to_json(response)

# Deserialize response
restored_response = ChannelSerializer.response_from_json(response_json)
assert restored_response.response_id == response.response_id

# Content hash for integrity checking
hash_val = ChannelSerializer.content_hash(request_json)
print(f"Content hash: {hash_val}")

# Full round-trip verification
rt_req = ChannelSerializer.round_trip_request(request)
rt_resp = ChannelSerializer.round_trip_response(response)
print(f"Round-trip request match: {rt_req.request_id}")
print(f"Round-trip response match: {rt_resp.response_id}")

# Jurisdiction serialization
juris = ChannelJurisdiction.for_channel(EvidenceChannel.COPILOT)
juris_json = ChannelSerializer.jurisdiction_to_json(juris)
restored_juris = ChannelSerializer.jurisdiction_from_json(juris_json)

# Configuration serialization
config = ChannelConfiguration.default_for(EvidenceChannel.COPILOT)
config_json = ChannelSerializer.config_to_json(config)
restored_config = ChannelSerializer.config_from_json(config_json)

# Pool status
pool_json = ChannelSerializer.pool_status_to_json({
    "copilot": "healthy",
    "solver": "healthy",
})
print(f"Pool status: {pool_json}")
\end{lstlisting}

% ─────────────────────────────────────────────────────────────────────
\section{Formal Properties}
\label{sec:formal}

We prove five formal properties of the $\CopChan$ and its
infrastructure.  All proofs are formalized in \LEAN{} in
\texttt{Paper76\_CopilotChannel.lean}.

\subsection{Jurisdiction Soundness}
\label{sec:formal:jurisdiction}

\begin{theorem}[Jurisdiction soundness]
\label{thm:jurisdiction}
For every request $q$ handled by the $\CopChan$ under jurisdiction
$J$, the request is admissible: $J \vdash q$.

Equivalently, no inadmissible request is processed:
\[
  \forall q.\; \CopChan.\CopHandle(q) \neq \bot
  \implies J \vdash q.
\]
\end{theorem}

\begin{proof}
The $\CopHandle$ method begins with an admissibility check:
$\AdmCheck(q.coordinate, q.proposition\_kind)$.  If the check
returns $\bot$ (inadmissible), the handler returns $\bot$ without
processing the request.  Therefore, any non-$\bot$ response
implies the check passed, which by Definition~\ref{def:admissibility}
means $J \vdash q$.

More formally, the control flow is:
\[
  \CopHandle(q) = \begin{cases}
    \bot & \text{if } J \nvdash q, \\
    \text{process}(q) & \text{if } J \vdash q.
  \end{cases}
\]
If $\CopHandle(q) \neq \bot$, we are in the second case, so $J \vdash q$.
\end{proof}

\subsection{Trust-Ceiling Idempotence}
\label{sec:formal:ceiling}

\begin{theorem}[Trust-ceiling idempotence]
\label{thm:ceiling-idemp}
For all trust tiers $\tau$ and ceiling $c$:
$\TrustCeil_c(\TrustCeil_c(\tau)) = \TrustCeil_c(\tau)$.
\end{theorem}

\begin{proof}
By Lemma~\ref{lem:ceil-idemp}.  Let $\tau' = \TrustCeil_c(\tau)$.
Since $\tau' \tleq c$, the enforcement function maps $\tau'$ to
itself.
\end{proof}

\begin{corollary}[Double enforcement is benign]
\label{cor:double-enforce}
The two-stage enforcement (channel + bridge) produces the same
result as single enforcement:
$\TrustCeil_c(\TrustCeil_c(\tau)) = \TrustCeil_c(\tau)$.
\end{corollary}

\begin{proof}
Direct consequence of Theorem~\ref{thm:ceiling-idemp}.
\end{proof}

\begin{theorem}[Ceiling enforcement is monotone]
\label{thm:ceil-mono}
For all $\tau_1 \tleq \tau_2$ and ceiling $c$:
$\TrustCeil_c(\tau_1) \tleq \TrustCeil_c(\tau_2)$.
\end{theorem}

\begin{proof}
Case analysis on the relationship between $\tau_1$, $\tau_2$, and $c$:
\begin{enumerate}
\item If $\tau_2 \tleq c$: then $\tau_1 \tleq \tau_2 \tleq c$, so
  $\TrustCeil_c(\tau_1) = \tau_1 \tleq \tau_2 = \TrustCeil_c(\tau_2)$.
\item If $\tau_1 \tleq c$ and $\tau_2 \not\tleq c$: then
  $\TrustCeil_c(\tau_1) = \tau_1 \tleq c = \TrustCeil_c(\tau_2)$.
\item If $\tau_1 \not\tleq c$: then since $\tau_1 \tleq \tau_2$
  and $c \tleq \tau_1$, we have $c \tleq \tau_2$, so
  $\TrustCeil_c(\tau_1) = c = \TrustCeil_c(\tau_2)$.
\end{enumerate}
In all cases, $\TrustCeil_c(\tau_1) \tleq \TrustCeil_c(\tau_2)$.
\end{proof}

\subsection{Rate-Limit Correctness}
\label{sec:formal:rate}

\begin{definition}[Rate-limited channel]
\label{def:rate-limited}
A channel with rate limit $\lambda > 0$ allows at most
$\lambda$ requests per second.  The \emph{rate-limit predicate}
for a sequence of request timestamps $t_1 < t_2 < \cdots < t_n$
within window $[T, T+1]$ is:
\[
  \mathrm{RateOK}(\lambda, t_1, \ldots, t_n) \iff
  |\{i : t_i \in [T, T+1]\}| \leq \lambda \quad \forall T.
\]
\end{definition}

\begin{theorem}[Rate-limit correctness]
\label{thm:rate-limit}
If the $\CopChan$ is configured with rate limit $\lambda > 0$,
then for any sequence of accepted requests with timestamps
$t_1, \ldots, t_n$:
$\mathrm{RateOK}(\lambda, t_1, \ldots, t_n)$.
\end{theorem}

\begin{proof}
The $\texttt{rate\_limit\_check}()$ method is invoked before
processing each request.  It maintains a sliding window counter
that tracks the number of requests in the current one-second
window.  If the counter equals $\lambda$, the request is rejected.
Therefore, at most $\lambda$ requests are accepted in any
one-second window, which is precisely $\mathrm{RateOK}$.
\end{proof}

\subsection{Round-Trip Fidelity}
\label{sec:formal:roundtrip}

\begin{theorem}[Serialization round-trip fidelity]
\label{thm:round-trip}
For every evidence request $r$ and evidence response $s$:
\begin{align*}
  \ChanSer.\texttt{from\_json}(\ChanSer.\texttt{to\_json}(r)) &= r, \\
  \ChanSer.\texttt{from\_json}(\ChanSer.\texttt{to\_json}(s)) &= s.
\end{align*}
\end{theorem}

\begin{proof}
The serializer uses a deterministic JSON encoding for all fields
of $\EvReq$ and $\EvResp$.  Each field type (string, integer,
float, enum, tuple, mapping) has a bijective encoding/decoding pair.
Since the composition of bijections is a bijection, the round-trip
preserves equality.

More precisely, for a request $r = (r_1, \ldots, r_k)$ with fields
$r_i$ of types $T_i$:
\[
  \texttt{from\_json}(\texttt{to\_json}(r))
  = (\texttt{dec}_{T_1}(\texttt{enc}_{T_1}(r_1)), \ldots,
     \texttt{dec}_{T_k}(\texttt{enc}_{T_k}(r_k)))
  = (r_1, \ldots, r_k) = r
\]
where each $\texttt{dec}_{T_i} \circ \texttt{enc}_{T_i} = \mathrm{id}_{T_i}$.
\end{proof}

\subsection{Monitor Completeness}
\label{sec:formal:monitor}

\begin{theorem}[Monitor completeness]
\label{thm:monitor}
If $\texttt{record\_request}$, $\texttt{record\_response}$, and
$\texttt{record\_failure}$ are called for every interaction, then
the monitor is complete (Definition~\ref{def:monitor-complete}).
\end{theorem}

\begin{proof}
By assumption, every request triggers a call to
$\texttt{record\_request}$, adding an entry to $L_R$.
Similarly for responses ($L_S$) and failures ($L_F$).
Therefore, for every interaction $I_k$, the corresponding log
contains a record, satisfying completeness.
\end{proof}

\begin{corollary}[Derived metrics are accurate]
\label{cor:metrics-accurate}
If the monitor is complete, then:
\begin{enumerate}
\item $\texttt{success\_rate}() = |L_S| / |L_R|$ (ignoring failures).
\item $\texttt{throughput}()$ correctly reflects the request rate.
\item $\texttt{latency\_percentiles}()$ are computed over all responses.
\end{enumerate}
\end{corollary}

\begin{proof}
Each derived metric is computed from the underlying logs.
Completeness ensures no interactions are missing, so the
metrics accurately reflect the true channel behavior.
\end{proof}

% ─────────────────────────────────────────────────────────────────────
\section{Discussion}
\label{sec:discussion}

\paragraph{Jurisdiction as access control.}
The jurisdiction mechanism serves as a form of access control
for the LLM: it determines which parts of the verification
task the LLM is allowed to contribute to.  A restrictive
jurisdiction limits the LLM to well-understood domains where
its suggestions are more likely to be useful.  A permissive
jurisdiction allows broader contribution but may increase
the rate of unhelpful suggestions.

\paragraph{Rate limiting vs.\ throughput.}
Rate limiting protects against resource exhaustion but reduces
throughput.  The optimal rate limit depends on the LLM provider's
capacity, the verification task's urgency, and the cost model.
In practice, a rate limit of 5--15 queries/second provides a
good balance.

\paragraph{Monitoring for degradation detection.}
The $\ChanMon$ enables proactive degradation detection: if the
success rate drops below a threshold or latency exceeds bounds,
the system can take corrective action (e.g., switching to a
fallback channel or increasing the timeout).

\paragraph{Serialization for reproducibility.}
Round-trip fidelity ensures that verification sessions can be
faithfully replayed.  This is essential for debugging, auditing,
and regression testing of the verification pipeline.

\paragraph{Design trade-offs.}
The two-stage trust enforcement (channel + bridge) adds
redundancy but prevents single-point-of-failure trust escalation.
The overhead is negligible (Theorem~\ref{thm:ceiling-idemp}
shows the second enforcement is a no-op if the first succeeds).

% ─────────────────────────────────────────────────────────────────────
\section{Related Work}
\label{sec:related}

\paragraph{LLM integration in verification.}
The integration of LLMs into formal verification has been explored
in several systems.  \citet{chen2021codex} evaluate code generation
capabilities; \citet{first2023baldur} use LLMs for proof generation.
Our work provides the channel infrastructure for trust-calibrated
integration.

\paragraph{Evidence channels and provenance.}
The concept of evidence channels extends prior work on proof
certificates~\cite{necula1997proof} and proof-carrying
code~\cite{appel2001foundational}.  Our contribution is the
formalization of jurisdiction, trust bounds, and infrastructure.

\paragraph{Rate limiting and API management.}
Rate limiting for API services has been studied in web
systems~\cite{nayak2024rate}.  We adapt these techniques to
the verification setting with formal correctness guarantees.

\paragraph{Serialization and reproducibility.}
Serialization with round-trip guarantees has been studied in
database systems~\cite{stonebraker1986design} and distributed
computing.  Our formalization applies these concepts to
evidence management.

\paragraph{Channel monitoring.}
Performance monitoring for distributed services is a mature
field~\cite{barham2004using}.  We specialize these concepts
to evidence channels with verification-specific metrics.

% ─────────────────────────────────────────────────────────────────────
\section{Conclusion}
\label{sec:conclusion}

We have presented a comprehensive study of the $\CopChan$
architecture, covering its implementation, jurisdiction mechanism,
configuration system, and infrastructure layer.  The five formal
properties established are:

\begin{enumerate}
\item \textbf{Jurisdiction soundness}
  (Theorem~\ref{thm:jurisdiction}): inadmissible requests are
  never processed.

\item \textbf{Trust-ceiling idempotence}
  (Theorem~\ref{thm:ceiling-idemp}): double enforcement is
  equivalent to single enforcement.

\item \textbf{Ceiling monotonicity}
  (Theorem~\ref{thm:ceil-mono}): lower input trust produces
  lower output trust.

\item \textbf{Rate-limit correctness}
  (Theorem~\ref{thm:rate-limit}): the configured rate is
  never exceeded.

\item \textbf{Round-trip fidelity}
  (Theorem~\ref{thm:round-trip}): serialization preserves
  message equality.

\item \textbf{Monitor completeness}
  (Theorem~\ref{thm:monitor}): all interactions are logged.
\end{enumerate}

The channel architecture provides a principled foundation for
integrating LLM evidence into formal verification with formal
guarantees about trust, jurisdiction, and infrastructure correctness.

% ─────────────────────────────────────────────────────────────────────
\begin{thebibliography}{25}

\bibitem{young-jugeo}
H.~Young.
\newblock Judgment Geometry: Semantic sites, descent, and the
  Copilot channel.
\newblock JuGeo Paper~0, 2024.

\bibitem{young-channels}
H.~Young.
\newblock The Copilot evidence channel.
\newblock JuGeo Paper~88, 2024.

\bibitem{young-trust}
H.~Young.
\newblock The trust algebra.
\newblock JuGeo Paper~4, 2024.

\bibitem{young-lifecycle}
H.~Young.
\newblock The Copilot lifecycle.
\newblock JuGeo Paper~93, 2024.

\bibitem{young-advisors}
H.~Young.
\newblock The advisor architecture.
\newblock JuGeo Paper~90, 2024.

\bibitem{chen2021codex}
M.~Chen, J.~Tworek, H.~Jun, et~al.
\newblock Evaluating large language models trained on code.
\newblock \emph{arXiv:2107.03374}, 2021.

\bibitem{first2023baldur}
E.~First, M.~Rabe, T.~Ringer, and Y.~Brun.
\newblock {Baldur}: Whole-proof generation and repair with large
  language models.
\newblock In \emph{Proc.\ ESEC/FSE}, 2023.

\bibitem{necula1997proof}
G.~C.~Necula.
\newblock Proof-carrying code.
\newblock In \emph{Proc.\ POPL}, 1997.

\bibitem{appel2001foundational}
A.~W.~Appel.
\newblock Foundational proof-carrying code.
\newblock In \emph{Proc.\ LICS}, 2001.

\bibitem{nayak2024rate}
R.~Nayak, S.~Patil, and K.~Joshi.
\newblock Rate limiting strategies for large-scale {API} services.
\newblock In \emph{Proc.\ ICWS}, 2024.

\bibitem{demoura2008z3}
L.~de~Moura and N.~Bj{\o}rner.
\newblock {Z3}: An efficient {SMT} solver.
\newblock In \emph{Proc.\ TACAS}, 2008.

\bibitem{de2015lean}
L.~de~Moura, S.~Kong, J.~Avigad, F.~van~Doorn, and
  J.~von~Raumer.
\newblock The {Lean} theorem prover (system description).
\newblock In \emph{Proc.\ CADE}, 2015.

\bibitem{stonebraker1986design}
M.~Stonebraker.
\newblock The design of the {POSTGRES} storage system.
\newblock In \emph{Proc.\ VLDB}, 1986.

\bibitem{barham2004using}
P.~Barham, A.~Donnelly, R.~Isaacs, and R.~Mortier.
\newblock Using {Magpie} for request extraction and workload
  modelling.
\newblock In \emph{Proc.\ OSDI}, 2004.

\bibitem{hoare1969axiomatic}
C.~A.~R.~Hoare.
\newblock An axiomatic basis for computer programming.
\newblock \emph{CACM}, 12(10):576--580, 1969.

\bibitem{dijkstra1975guarded}
E.~W.~Dijkstra.
\newblock Guarded commands, nondeterminacy and formal derivation
  of programs.
\newblock \emph{CACM}, 18(8):453--457, 1975.

\bibitem{leino2010dafny}
K.~R.~M.~Leino.
\newblock {Dafny}: An automatic program verifier for functional
  correctness.
\newblock In \emph{Proc.\ LPAR}, 2010.

\bibitem{ouyang2022training}
L.~Ouyang, J.~Wu, X.~Jiang, et~al.
\newblock Training language models to follow instructions with
  human feedback.
\newblock In \emph{Proc.\ NeurIPS}, 2022.

\bibitem{li2022competition}
Y.~Li, D.~Choi, J.~Chung, et~al.
\newblock Competition-level code generation with {AlphaCode}.
\newblock \emph{Science}, 378(6624):1092--1097, 2022.

\bibitem{coq2024manual}
The Coq Development Team.
\newblock \emph{The Coq Reference Manual}, 2024.

\bibitem{floyd1967assigning}
R.~W.~Floyd.
\newblock Assigning meanings to programs.
\newblock In \emph{Proc.\ Symp.\ Applied Mathematics}, 1967.

\end{thebibliography}

\appendix
\section{Additional Proofs and Properties}
\label{app:proofs}

\subsection{Trust Lattice Completeness}
\label{app:lattice}

\begin{lemma}[Trust lattice is a bounded lattice]
\label{lem:bounded-lattice}
The trust lattice $(\mathcal{T}, \tleq)$ with elements
$\Tcontra \tleq \Tunver \tleq \Tcopilot \tleq \Toracle
\tleq \Truntime \tleq \Tsolver \tleq \Tproof$
is a bounded lattice with bottom $\Tcontra$ and top $\Tproof$.
\end{lemma}

\begin{proof}
The set $\mathcal{T}$ is totally ordered by $\tleq$, so it is
automatically a lattice with meet $\tjoin = \min$ and join
$\tatten = \max$.  The bottom element is $\Tcontra$ (minimum
under $\tleq$) and the top element is $\Tproof$ (maximum).
\end{proof}

\begin{lemma}[Ceiling is a meet-preserving map]
\label{lem:ceil-meet}
For all $\tau_1, \tau_2$ and ceiling $c$:
$\TrustCeil_c(\tau_1 \tjoin \tau_2) = \TrustCeil_c(\tau_1) \tjoin \TrustCeil_c(\tau_2)$.
\end{lemma}

\begin{proof}
Let $m = \tau_1 \tjoin \tau_2 = \min(\tau_1, \tau_2)$.
We need to show $\min(m, c) = \min(\min(\tau_1, c), \min(\tau_2, c))$.
Both sides equal $\min(\tau_1, \tau_2, c)$ by associativity and
commutativity of $\min$.
\end{proof}

\subsection{Jurisdiction Algebraic Properties}
\label{app:jurisdiction}

\begin{definition}[Jurisdiction intersection]
\label{def:juris-intersect}
For jurisdictions $J_1 = (D_1, P_1, K_1, \tau_1, r_1, X_1)$ and
$J_2 = (D_2, P_2, K_2, \tau_2, r_2, X_2)$, their intersection is:
\[
  J_1 \cap J_2 = (D_1 \cap D_2, P_1 \cap P_2, K_1 \cap K_2,
  \min(\tau_1, \tau_2), r_1 \vee r_2, X_1 \cup X_2).
\]
\end{definition}

\begin{lemma}[Intersection preserves admissibility]
\label{lem:intersect-admis}
If $J_1 \cap J_2 \vdash q$ then $J_1 \vdash q$ and $J_2 \vdash q$.
\end{lemma}

\begin{proof}
Admissibility under $J_1 \cap J_2$ requires the coordinate to
match a pattern in $P_1 \cap P_2 \subseteq P_1$ (so $J_1$'s
pattern check passes) and similarly for $P_2$.  The proposition
kind must be in $K_1 \cap K_2 \subseteq K_1$ and $K_2$.  The
query family must not be in $X_1 \cup X_2 \supseteq X_1$ and
$X_2$, which is a stronger condition.  However, the statement
says $J_1 \cap J_2 \vdash q$, meaning the query family is not
in $X_1 \cup X_2$, so it is not in $X_1$ alone nor in $X_2$ alone.
Therefore $J_1 \vdash q$ and $J_2 \vdash q$.
\end{proof}

\begin{lemma}[Jurisdiction intersection is commutative]
\label{lem:juris-comm}
$J_1 \cap J_2 = J_2 \cap J_1$.
\end{lemma}

\begin{proof}
Each component operation ($\cap$, $\min$, $\vee$, $\cup$) is
commutative.
\end{proof}

\begin{lemma}[Jurisdiction intersection is associative]
\label{lem:juris-assoc}
$(J_1 \cap J_2) \cap J_3 = J_1 \cap (J_2 \cap J_3)$.
\end{lemma}

\begin{proof}
Each component operation is associative.
\end{proof}

\subsection{Rate Limit Window Properties}
\label{app:rate}

\begin{lemma}[Sliding window is monotone]
\label{lem:window-mono}
If $\mathrm{RateOK}(\lambda, t_1, \ldots, t_n)$ and $t_{n+1}$
falls outside all existing windows that are at capacity, then
$\mathrm{RateOK}(\lambda, t_1, \ldots, t_{n+1})$.
\end{lemma}

\begin{proof}
Adding $t_{n+1}$ increments the count in the window
$[t_{n+1}-1, t_{n+1}]$ by one.  Since $t_{n+1}$ was chosen
outside capacity windows, this count was at most $\lambda - 1$
before the addition, so it is at most $\lambda$ after.
\end{proof}

\begin{lemma}[Rate limit rejection is correct]
\label{lem:rate-reject}
If a request at time $t$ is rejected by the rate limiter, then
accepting it would violate $\mathrm{RateOK}$.
\end{lemma}

\begin{proof}
The rate limiter rejects a request at time $t$ only if the
window $[t-1, t]$ already contains $\lambda$ requests.  Adding
the request would bring the count to $\lambda + 1 > \lambda$,
violating $\mathrm{RateOK}$.
\end{proof}

\subsection{Monitor Log Properties}
\label{app:monitor}

\begin{lemma}[Monitor logs are append-only]
\label{lem:log-append}
Each call to $\texttt{record\_request}$, $\texttt{record\_response}$,
or $\texttt{record\_failure}$ strictly increases the size of the
corresponding log: $|L| \to |L| + 1$.
\end{lemma}

\begin{proof}
The record methods append to the log without removing existing
entries.
\end{proof}

\begin{lemma}[Log ordering preserves temporal order]
\label{lem:log-order}
Entries in each log are ordered by timestamp:
$e_i.\mathrm{timestamp} \leq e_{i+1}.\mathrm{timestamp}$
for consecutive entries $e_i, e_{i+1}$.
\end{lemma}

\begin{proof}
Each record call uses the current time as the timestamp.  Since
calls are sequential and time is monotonically increasing, the
ordering is preserved.
\end{proof}

\end{document}
